% FORMALIZACION_NUEVA of the pre-existing homogeneous realization.
% Priority of the integrated formula: historical provenance pending comparison.
% Preserves the equation and domain of the bounce chapter; does not select s_0.
\section{Explicit homogeneous solution and causal completeness}
\label{vii:sec:solucion-homogenea}
Determining the amplitude permits a complete realization of the bounce
to be developed in addition to the local argument. The spatially flat
chart of \eqref{vii:eq:dinamica-homogenea} is retained, in units
\(c=1\), with \(\rho_0,s_0,G>0\), and the material sector is specialized
to dust, \(w=0\), with \(\Lambda_{\mathrm{eff}}=0\).
The coefficient \(s_0\) is obtained from the material state in the
corresponding realization; the following integration uses this
coefficient and preserves its domain.

\begin{theorem}[Exact integration of the dust bounce]
\label{vii:thm:rebote-polvo}
Set
\[
 a_{\rm b}^3=\frac{s_0}{\rho_0},\qquad
 \tau^2=\frac{s_0}{6\pi G\rho_0^2},\qquad u=t-t_{\rm b}.
\]
The solution of the homogeneous system with
\(a(t_{\rm b})=a_{\rm b}\) and \(H(t_{\rm b})=0\) is, for every
\(t\in\mathbb R\),
\begin{equation}
 a(t)=a_{\rm b}\left(1+\frac{u^2}{\tau^2}\right)^{1/3},
 \qquad H(t)=\frac{2u}{3(\tau^2+u^2)},\qquad
 \dot H(t)=\frac{2(\tau^2-u^2)}{3(\tau^2+u^2)^2}.
 \label{vii:eq:solucion-polvo}
\end{equation}
It has a unique minimum at \(t_{\rm b}\); it is contracting before
and expanding afterwards. The evolution is smooth and satisfies
\(a(t)\geq a_{\rm b}>0\).
\end{theorem}
\begin{proof}
Let \(y=a^3\). The Friedmann equation gives
\[
 \dot y^{\,2}=24\pi G(\rho_0y-s_0).
\]
The function \(y=s_0/\rho_0+6\pi G\rho_0u^2\) satisfies this
identity and the initial conditions. Its derivatives give the
three expressions in \eqref{vii:eq:solucion-polvo}. To verify the
evolution equation as well, write
\[
 \rho_{\rm b}=\frac{\rho_0^2}{s_0},\qquad
 x=\frac{u}{\tau}.
\]
Then
\[
 \rho_{\rm m}=\frac{\rho_{\rm b}}{1+x^2},\qquad
 \rho_{\rm s}=\frac{\rho_{\rm b}}{(1+x^2)^2},\qquad
 4\pi G\rho_{\rm b}=\frac{2}{3\tau^2}.
\]
Substitution into \(\dot H=-4\pi G(\rho_{\rm m}-2\rho_{\rm s})\)
gives exactly the third formula. Thus the complete system
\((\dot a,\dot H)\) is satisfied, and its vector field is smooth
for \(a>0\). Local uniqueness, prolonged along the explicit
solution, determines the evolution from these initial conditions.
The denominator is positive for every \(t\), and the statements
about the minimum and sign follow from \(H\).
\end{proof}

\begin{figure}[htbp]
\centering
\begin{tikzpicture}[x=1.45cm,y=1.15cm,>=Latex]
\draw[->] (-3.25,0)--(3.3,0) node[right] {\(\displaystyle u/\tau\)};
\draw[->] (0,0)--(0,2.8) node[above] {\(\displaystyle a/a_{\rm b}\)};
\draw[densely dashed,gray] (-3.1,1)--(3.1,1);
\draw[very thick,ink,domain=-3:3,samples=121,smooth]
 plot (\x,{(1+\x*\x)^(1/3)});
\fill[accent] (0,1) circle (1.6pt);
\node[anchor=north east] at (-.08,.98) {\(1\)};
\node[ink] at (-2,2.65) {Contraction};
\node[ink] at (2,2.65) {Expansion};
\foreach \t in {-3,-2,-1,1,2,3}
{\draw (\t,.045)--(\t,-.045) node[below] {\(\t\)};}
\end{tikzpicture}
\caption{Dimensionless profile of the proved solution:
\(a/a_{\rm b}=(1+(u/\tau)^2)^{1/3}\).
The parameters \(a_{\rm b}\) and \(\tau\) are calculated from
\(\rho_0,s_0,G\); the diagram represents the entire rescaled collection.}
\label{vii:fig:rebote-polvo}
\end{figure}

\subsection{Finite curvature and extension of geodesics}
The metric realized in this chart is
\(g=-dt^2+a(t)^2\,d\mathbf x^2\) on
\(\mathbb R\times\mathbb R^3\). With the convention in which
the flat-model scalar curvature is \(R=6(\dot H+2H^2)\), the solution gives
\[
 R(t)=\frac{4\tau^2+\frac43u^2}{(\tau^2+u^2)^2}.
\]
In particular, \(R(t_{\rm b})=4/\tau^2\) is finite.
In an orthonormal frame, the Riemann components depend on
\(H^2\) and \(\dot H+H^2\), which are bounded functions on
\(\mathbb R\). The same holds for polynomial contractions of
these components; for example,
\[
 R_{\mu\nu\rho\sigma}R^{\mu\nu\rho\sigma}
 =12\bigl[(\dot H+H^2)^2+H^4\bigr].
\]

\begin{proposition}[Completeness of causal geodesics]
\label{vii:prop:completitud-causal-polvo}
The preceding metric realization is geodesically complete
for the timelike and null geodesics of its Levi--Civita
connection, in both directions.
\end{proposition}
\begin{proof}
Spatial translations provide the constants
\(p_i=a^2\,dx^i/d\lambda\). Causal normalization gives
\[
 \left|\frac{dt}{d\lambda}\right|
 =\sqrt{\epsilon+\frac{|\mathbf p|^2}{a(t)^2}},
 \qquad \epsilon=1\ \text{in the timelike case},\quad
 \epsilon=0\ \text{in the null case}.
\]
A nonconstant null geodesic has \(|\mathbf p|>0\).
Since \(a\geq a_{\rm b}\), in the timelike case
\[
 \left|\frac{d\lambda}{dt}\right|
 \geq\left(1+\frac{|\mathbf p|^2}{a_{\rm b}^2}\right)^{-1/2}>0,
\]
and in the null case
\(\left|d\lambda/dt\right|=a/|\mathbf p|
 \geq a_{\rm b}/|\mathbf p|>0\).
Thus reaching \(t=\pm\infty\) requires infinite affine length.
On a compact time interval, \(a\) and its derivatives are
smooth and \(a\) is bounded away from zero. The explicit equations
above and \(dx^i/d\lambda=p_i/a^2\) allow the geodesic to be
continued through its finite endpoints. Both possibilities
for prolongation are covered.
\end{proof}

The completeness result belongs to this homogeneous dust realization,
with its conditions on space, matter and amplitude. Local persistence
under threshold perturbations preserves the bounds of
Theorem~\ref{vii:thm:persistencia}; this local property and global
completeness of the explicit solution are distinct conclusions.
The minimum scale follows from the relation between material density
and torsional response. Memory thus enters a verifiable geometric
evolution, with a regular realization of contraction and expansion.
